% Values: PRECISE, FAITHFUL, TRACTABLE, ILLUMINATING, ECONOMICAL
% Evidence planned:
% - PRECISE: a rule (or an option instance of one) for every production of the
%   bare Core in the book plan's "Bare Constructs" table: 45 rules in 12 groups;
%   every object and operation defined in the sections before the rules.
% - FAITHFUL: classes and admissible types exactly as the survey's tables (A19,
%   A36); rows with unique labels (A23); finite types (A22); explicit variables
%   scoped over nested annotations (A21) and never escaping (A20); groups
%   monomorphic within themselves (A24); no value restriction; Int constants of
%   class number, Int patterns of type Int; unrestricted equality (A35);
%   exhaustiveness, redundancy by value (A31), inhabitants possibly infinite
%   (A97); principal derivations (A95), on which legality is judged (A96).
% - TRACTABLE: the property that Adm is monotone in the class order and closed
%   under substitutions that respect classes, so that instances compose.
% - ILLUMINATING: comments on instantiation, the freedom of integer constants,
%   record update, the two rules for tyvar, the escape condition, the group rule
%   and the absence of a value restriction, constructor status in patterns.
% - ECONOMICAL: options merge annotated and unannotated value bindings and
%   optional rows; derived forms (if, &&, ||, negation, lists, accessors, access,
%   multi-pattern lambdas, function definitions) have no rules here.

\chapter{Static Semantics for the Core}\label{ch:static-core}

This chapter defines the semantic objects of the static semantics of the Core
and the rules by which Core phrases elaborate. The simple objects are type
variables, each with a class, row variables, type names and identifier status
descriptors. From them we build types, record types, type functions, type
schemes, environments and contexts, and on these we define substitution,
instantiation and closure. The rules infer sentences $\jdg{C}{\mathit{phrase}}{A}$,
in which $C$ is a context and $A$ is a type, a row, a value environment, or a
value environment paired with a type. A Core phrase is legal when it satisfies
the syntactic restrictions of Chapter~\ref{ch:syntax-core}, elaborates by the
rules of Section~\ref{sec:static-core-rules}, and its principal derivation
(Section~\ref{sec:static-core-principal}) meets the further restrictions of
Section~\ref{sec:static-core-further}. % A96

\section{Simple Objects}\label{sec:static-core-simple}

All semantic objects of the static semantics are built from identifiers
(Chapter~\ref{ch:syntax-core}), labels among them, and four further kinds of
simple object: type variables, row variables, type names and identifier status
descriptors, together with the classes that type variables carry. Type names
are the values taken by type constructors; they are distinct from type
variables and from type constructors. The simple object classes, and the
metavariables ranging over them, are shown in Figure~\ref{fig:static-core-simple}.
We include $\TyVar$ in the figure to show that $\alpha$ ranges over it.

\begin{figure}[ht]
\centering
\begin{objtable}
  \objrow{\alpha\ \text{or}\ \tyvar}{\TyVar}{type variables}
  \objrow{\epsilon}{\RowVar}{row variables}
  \objrow{t}{\TyName}{type names}
  \objrow{\kappa}{\Class}{classes}
  \objrow{\is}{\IdStatus = \{\isc, \isv\}}{identifier status descriptors}
\end{objtable}
\caption{Simple Semantic Objects}\label{fig:static-core-simple}
\end{figure}

The classes are
\[
  \Class = \{\kany, \knum, \kcomp, \kapp, \kcapp\}.
\]
Each $\alpha \in \TyVar$ has a class $\clsof(\alpha) \in \Class$, and there
are infinitely many type variables of each class. A row variable has no class.
Each $t \in \TyName$ has an arity $k \ge 0$; we write $\TyName^{(k)}$ for the
type names of arity $k$.

A type variable $\tyvar$ of the syntax names one type variable, also written
$\tyvar$, whose class is fixed by the beginning of its name: $\knum$ if the
name begins with \kw{number}, $\kcomp$ if it begins with \kw{comparable},
$\kapp$ if it begins with \kw{appendable}, $\kcapp$ if it begins with
\kw{compappend}, and $\kany$ otherwise. % A19
When that class is $\kany$, $\tyvar$ also names one row variable; a type
variable of the syntax of any other class names no row variable. % A36
We write $\tyvar = \alpha$ and $\tyvar = \epsilon$ when $\alpha$, respectively
$\epsilon$, is the variable named by $\tyvar$. No other type or row variable is
named by a type variable of the syntax.

A label and a value variable may be spelt alike. For a label $\lab$ we write
$\var = \lab$ when $\var$ is spelt as $\lab$.

The rules mention a few type names directly. They are shown in
Figure~\ref{fig:static-core-builtin}; they belong to the initial basis
(Appendix~\ref{ch:appendix-static-basis}), and a rule that names one of them
means that type name, whatever a program binds to the type constructor of the
same spelling.

\begin{figure}[ht]
\centering
\begin{objtable}
  \objrow{\tInt,\ \tFloat}{\TyName^{(0)}}{integers, floats}
  \objrow{\tChar,\ \tString}{\TyName^{(0)}}{characters, strings}
  \objrow{\tunit}{\TyName^{(0)}}{unit}
  \objrow{\tList}{\TyName^{(1)}}{lists}
  \objrow{\tTuple{2}}{\TyName^{(2)}}{pairs}
  \objrow{\tTuple{3}}{\TyName^{(3)}}{triples}
\end{objtable}
\caption{Built-in Type Names}\label{fig:static-core-builtin}
\end{figure}

\section{Compound Objects}\label{sec:static-core-compound}

When $A$ and $B$ are sets, $\Fin A$ denotes the set of finite subsets of $A$,
and $A \finmap B$ denotes the set of \emph{finite maps} (partial functions with
finite domain) from $A$ to $B$. The domain and range of a finite map $f$ are
denoted $\Dom f$ and $\Ran f$. A finite map is written explicitly as
$\{a_1 \mapsto b_1, \cdots, a_k \mapsto b_k\}$, $k \ge 0$; in particular the
empty map is $\emptymap$. We write $\{x \mapsto e \,;\, \phi\}$ for the finite
map $f$ whose domain is the set of $x$ satisfying the condition $\phi$ and whose
value on that domain is $f(x) = e$. A finite map is regarded as the set of its
pairs, so $f \subseteq g$ means that $\Dom f \subseteq \Dom g$ and that $g$
agrees with $f$ on $\Dom f$.

When $f$ and $g$ are finite maps, the map $f + g$, called $f$
\emph{modified} by $g$, is the finite map with domain $\Dom f \cup \Dom g$ and
values
\[
  (f + g)(a) = \text{if } a \in \Dom g \text{ then } g(a) \text{ else } f(a).
\]

The compound objects of the static semantics of the Core are shown in
Figure~\ref{fig:static-core-compound}. We take $\cup$ to mean disjoint union
over semantic object classes, and we understand all the defined object classes
to be disjoint. The classes $\Type$, $\Row$, $\RecType$, $\FunType$ and
$\ConsType$ are the smallest sets satisfying their equations; hence every type
is finite, and no type occurs within itself. % A22
A sequence $(\tau_1, \cdots, \tau_k)$ of types is also written $\tau^{(k)}$, and
a sequence $(\tyrowvar_1, \cdots, \tyrowvar_k)$ of type and row variables
$\tyrowvar^{(k)}$; the metavariable $\tyrowvar$ ranges over $\TyVar \cup
\RowVar$. In a context, $T \in \Fin(\TyName)$ and $U \in \Fin(\TyVar \cup
\RowVar)$.

\begin{figure}[ht]
\centering
\footnotesize
\begin{objtable}
  \objdef{\tau}{\Type}{\TyVar \cup \RecType \cup \FunType \cup \ConsType}
  \objdef{\varrho}{\Row}{\cls{Lab} \finmap \Type}
  \objdef{\rec{\varrho}\ \text{or}\ \orec{\varrho}{\epsilon}}{\RecType}{\Row \times (\RowVar \cup \{\emptyrow\})}
  \objdef{\tau \rightarrow \tau'}{\FunType}{\Type \times \Type}
  \objdef{}{\ConsType}{\textstyle\bigcup_{k \ge 0} \ConsType^{(k)}}
  \objdef{\tau^{(k)} t}{\ConsType^{(k)}}{\Type^k \times \TyName^{(k)}}
  \objdef{\theta\ \text{or}\ \Lambda\tyrowvar^{(k)}.\tau}{\TypeFcn}{\textstyle\bigcup_{k \ge 0} (\TyVar \cup \RowVar)^k \times \Type}
  \objdef{\sigma\ \text{or}\ \forall\tyrowvar^{(k)}.\tau}{\TypeScheme}{\textstyle\bigcup_{k \ge 0} (\TyVar \cup \RowVar)^k \times \Type}
  \objdef{(\theta, \VE)}{\TyStr}{\TypeFcn \times \ValEnv}
  \objdef{\TE}{\TyEnv}{\cls{TyCon} \finmap \TyStr}
  \objdef{\VE}{\ValEnv}{(\cls{Var} \cup \cls{Con}) \finmap \TypeScheme \times \IdStatus}
  \objdef{\QE}{\QualEnv}{\cls{ModId} \finmap \Env}
  \objdef{E\ \text{or}\ \QE, \TE, \VE}{\Env}{\QualEnv \times \TyEnv \times \ValEnv}
  \objdef{C\ \text{or}\ T, U, E}{\Context}{\Fin(\TyName) \times \Fin(\TyVar \cup \RowVar)}
  & & & \quad{}\times \Env & \\
\end{objtable}
\caption{Compound Semantic Objects}\label{fig:static-core-compound}
\end{figure}

A record type is a pair of a row and either a row variable or $\emptyrow$. The
record type $(\varrho, \emptyrow)$, written $\rec{\varrho}$, is \emph{closed}:
its fields are exactly those of $\varrho$. The record type $(\varrho,
\epsilon)$, written $\orec{\varrho}{\epsilon}$, is \emph{open}: its fields are
those of $\varrho$ together with those that $\epsilon$ stands for, none of which
has a label in $\Dom\varrho$ (Section~\ref{sec:static-core-types}). % A23
The unit type is the constructed type $()\tunit$, and the type of $k$-tuples
with components $\tau_1, \cdots, \tau_k$ is $(\tau_1, \cdots,
\tau_k)\tTuple{k}$. When $t$ has arity 0 we write $t$ for the type $()t$; thus
$\tInt$, $\tunit$ and so on are types.

Note that $\Lambda$ and $\forall$ bind type and row variables. For any
semantic object $A$, $\tynames A$ and $\tyvars A$ denote respectively the set
of type names and the set of type and row variables occurring free in $A$. For
a context we take $\tyvars C = \tyvars(E \of C) \cup (U \of C)$. For a type
expression $\ty$, $\tyvars \ty$ denotes the set of type variables of the
syntax occurring in $\ty$.

A value environment maps value variables and value constructors to a pair of a
type scheme and an identifier status. If $\VE(\var) = (\sigma, \is)$ or
$\VE(\con) = (\sigma, \is)$, we say that $\var$, respectively $\con$, \emph{has
status} $\is$ in $\VE$. Value variables have status $\isv$. The constructors of
custom types have status $\isc$, and only they may occur in patterns; the
constructor that a record type alias defines has status $\isv$
(Chapter~\ref{ch:static-modules}). We take $\cls{Con}$ to contain, besides
the constructors spelt by identifiers, the two list constructor names \kw{[]}
and \kw{::}, which no identifier spells; the initial basis binds them with
status $\isc$ to $\forall\alpha.\alpha\,\tList$ and $\forall\alpha.\alpha
\rightarrow \alpha\,\tList \rightarrow \alpha\,\tList$, where $\clsof(\alpha) =
\kany$ (Appendix~\ref{ch:appendix-static-basis}).

A qualified environment maps each qualifier to the environment of the names it
qualifies; every environment in the range of a qualified environment has the
empty qualified environment. In a context $C = T, U, E$, $T$ is the set of type
names generated so far, and $U$ is the set of explicit type and row variables in
scope (Section~\ref{sec:static-core-explicit}).

\section{Projection, Injection and Modification}\label{sec:static-core-projection}

\textbf{Projection:} We often select a component of a tuple of objects, and we
rely on metavariable names to say which component. For instance ``$\VE \of
E$'' means the value environment component of $E$, and, for a record type
$\tau$, ``$\varrho \of \tau$'' means its row.

When a tuple contains a finite map we apply the tuple to an argument, relying
on the syntactic class of the argument to select the map. Thus $C(\var)$ means
$(\VE \of (E \of C))(\var)$, and likewise $C(\con)$; $C(\tycon)$
means $(\TE \of (E \of C))(\tycon)$. A long identifier is looked up through the
qualified environment: if $\longvar = \modid.\var$ then $E(\longvar)$ means
$(\VE \of (\QE \of E)(\modid))(\var)$, and likewise for $\longcon$ and
$\longtycon$; an unqualified $\longvar$ is a $\var$. $C(\longvar)$ means $(E
\of C)(\longvar)$. The identifiers that the equivalent forms of derived forms
introduce (Appendix~\ref{ch:appendix-derived-forms}) are looked up in the same
way. Because only an import of the module \kw{Basics} provides the qualifier
\kw{Basics} (Section~\ref{sec:syntax-modules-restrictions}), \kw{Basics.True},
\kw{Basics.False} and \kw{Basics.negate} always denote the bindings of the
initial basis (Appendix~\ref{ch:appendix-static-basis}); the constructor
\kw{::} is bound in the initial environment $E_0$ of that appendix, from which
every module is elaborated. % audit X18, V3; A89

Operators are not in the domain of an environment, and $C(\opv)$ is not
defined. Since no program declares an operator, the operators are fixed by the
initial basis: Appendix~\ref{ch:appendix-static-basis} gives the \emph{static
operator table}
\[
  \OTstat : \cls{Op} \finmap \TypeScheme \times \IdStatus ,
\]
which maps each operator of Figure~\ref{fig:core-fixities} to the type scheme
and status of the identifier that the operator stands for; in particular
$\OTstat(\kw{::})$ is the type scheme of the list constructor \kw{::}, with
status $\isc$. The operators of library modules outside the initial basis are
outside the definition. % A16; audit X1

\textbf{Injection:} Components may be injected into tuple classes; for example,
``$\VE$ in $\Env$'' means the environment $(\emptymap, \emptymap, \VE)$.

\textbf{Modification:} The modification of one map by another has been defined
above. It is extended to environments component by component, $E + E'$, and
empty components are left implicit: $E + \VE$ means $E + (\emptymap,
\emptymap, \VE)$. For set components modification means union, so that $C +
(T, U, \VE)$ means
\[
  \big((T \of C) \cup T,\ (U \of C) \cup U,\ (E \of C) + \VE\big),
\]
and $C + \VE$ and $C + U$ leave the other components unchanged. Finally,
$C \oplus \TE$ means $C + (\tynames \TE, \TE)$: the type environment is added
and $T$ is extended with its type names.

\section{Classes and Admissible Types}\label{sec:static-core-classes}

The classes are ordered by $\leq$, the least partial order on $\Class$ such
that
\[
\begin{array}{@{}c@{}}
  \knum \leq \kcomp, \qquad \kcapp \leq \kcomp,\\
  \kcapp \leq \kapp, \qquad \kappa \leq \kany\ \text{for every}\ \kappa \in \Class.
\end{array}
\]
For each class $\kappa$, the set $\Adm{\kappa} \subseteq \Type$ of types
\emph{admissible} for $\kappa$ is defined in Figure~\ref{fig:static-core-admissible};
$\Adm{\kcomp}$ is the smallest set satisfying its equation. A type variable of
class $\kappa$ may stand only for a type in $\Adm{\kappa}$
(Section~\ref{sec:static-core-types}). A type variable belongs to
$\Adm{\kappa}$ exactly when its class is at most $\kappa$; no record type and no
function type is admissible for any class but $\kany$.

\begin{figure}[ht]
\centering
\[
\begin{array}{@{}r@{\ }c@{\ }l@{}}
  \Adm{\kany} & = & \Type\\[2pt]
  \Adm{\knum} & = & \{\tInt, \tFloat\} \cup \{\alpha \,;\, \clsof(\alpha) = \knum\}\\[2pt]
  \Adm{\kcomp} & = & \{\tInt, \tFloat, \tChar, \tString\}\\
  & \cup & \{\alpha \,;\, \clsof(\alpha) \leq \kcomp\}\\
  & \cup & \{\tau\,\tList \,;\, \tau \in \Adm{\kcomp}\}\\
  & \cup & \{(\tau_1, \tau_2)\tTuple{2} \,;\, \tau_1, \tau_2 \in \Adm{\kcomp}\}\\
  & \cup & \{(\tau_1, \tau_2, \tau_3)\tTuple{3} \,;\, \tau_1, \tau_2, \tau_3 \in \Adm{\kcomp}\}\\[2pt]
  \Adm{\kapp} & = & \{\tString\} \cup \{\tau\,\tList \,;\, \tau \in \Type\}\\
  & \cup & \{\alpha \,;\, \clsof(\alpha) \leq \kapp\}\\[2pt]
  \Adm{\kcapp} & = & \{\tString\} \cup \{\tau\,\tList \,;\, \tau \in \Adm{\kcomp}\}\\
  & \cup & \{\alpha \,;\, \clsof(\alpha) = \kcapp\}
\end{array}
\]
\caption{Admissible Types}\label{fig:static-core-admissible}
\end{figure}

There is no class of types that admit equality. The initial basis gives the
equality operators the type scheme $\forall\alpha.\alpha \rightarrow \alpha
\rightarrow \tBool$ with $\clsof(\alpha) = \kany$
(Appendix~\ref{ch:appendix-static-basis}), so equality elaborates at every
type, function types included. % A35

\section{Types and Type Functions}\label{sec:static-core-types}

\textbf{Substitution.} Let $\tyrowvar^{(k)}$ be distinct type and row
variables, and $\tau^{(k)}$ types such that $\tau_i \in \RecType$ whenever
$\tyrowvar_i \in \RowVar$. For a type $\tau$, the type $\tau\{\tau^{(k)} /
\tyrowvar^{(k)}\}$ is the result of replacing each $\tyrowvar_i$ by $\tau_i$
throughout $\tau$. Writing $\tau\{\cdot\}$ for $\tau\{\tau^{(k)} /
\tyrowvar^{(k)}\}$, and $\varrho\{\cdot\}$ for $\{\lab \mapsto
\varrho(\lab)\{\cdot\} \,;\, \lab \in \Dom\varrho\}$, it is defined by
\[
\begin{array}{@{}r@{\ =\ }l@{\qquad}l@{}}
  \alpha\{\cdot\} & \tau_i & \text{if}\ \alpha = \tyrowvar_i\\
  \alpha\{\cdot\} & \alpha & \text{if}\ \alpha \notin \{\tyrowvar_1, \cdots, \tyrowvar_k\}\\
  \rec{\varrho}\{\cdot\} & \rec{\varrho\{\cdot\}} &\\
  \orec{\varrho}{\epsilon}\{\cdot\} & \rec{\varrho\{\cdot\} + \varrho'} & \text{if}\ \epsilon = \tyrowvar_i\ \text{and}\ \tau_i = \rec{\varrho'}\\
  \orec{\varrho}{\epsilon}\{\cdot\} & \orec{\varrho\{\cdot\} + \varrho'}{\epsilon'} & \text{if}\ \epsilon = \tyrowvar_i\ \text{and}\ \tau_i = \orec{\varrho'}{\epsilon'}\\
  \orec{\varrho}{\epsilon}\{\cdot\} & \orec{\varrho\{\cdot\}}{\epsilon} & \text{if}\ \epsilon \notin \{\tyrowvar_1, \cdots, \tyrowvar_k\}\\
  (\tau \rightarrow \tau')\{\cdot\} & \tau\{\cdot\} \rightarrow \tau'\{\cdot\} &\\
  (\tau'_1, \cdots, \tau'_l)t\{\cdot\} & (\tau'_1\{\cdot\}, \cdots, \tau'_l\{\cdot\})t &
\end{array}
\]
The result is defined only if, wherever a row variable is replaced, $\Dom\varrho
\cap \Dom\varrho' = \emptyset$. Thus a row variable can never give a record
type a second field with a label it already has. % A23
The substitution \emph{respects classes} if $\tau_i \in
\Adm{\clsof(\tyrowvar_i)}$ for each $\tyrowvar_i \in \TyVar$. Substitution is
extended to rows, as above, and to other objects in the obvious way, renaming
bound variables where necessary.

The rules rely on the following property, which may easily be demonstrated.

\begin{property}
If $\kappa \leq \kappa'$ then $\Adm{\kappa} \subseteq \Adm{\kappa'}$.
If $\tau \in \Adm{\kappa}$, and $\tau\{\tau^{(k)}/\tyrowvar^{(k)}\}$ is defined
and respects classes, then $\tau\{\tau^{(k)}/\tyrowvar^{(k)}\} \in
\Adm{\kappa}$.
\end{property}

\textbf{Type functions.} A type function $\theta = \Lambda\tyrowvar^{(k)}.\tau$
has arity $k$; the bound variables must be distinct. Two type functions are
equal if they differ only in their choice of bound variables, a type variable
being renamed to a type variable and a row variable to a row variable. The class
of a bound type variable of a type function has no significance.
% DECIDED: the arguments of a type function (a type alias or a type name) are not
% required to be admissible for the classes of its parameters, as in the
% implementation, which substitutes alias arguments directly
% (compiler/src/Compiler/Type/Instantiate.elm).
If $t$ has arity $k$, then we write $t$ for the type function
$\Lambda\alpha^{(k)}.\alpha^{(k)}t$, where the $\alpha_i$ are distinct type
variables; thus $\TyName \subseteq \TypeFcn$.

The application of $\theta = \Lambda\tyrowvar^{(k)}.\tau$ to types
$\tau^{(k)}$ is
\[
  \theta(\tau^{(k)}) = \tau\{\tau^{(k)}/\tyrowvar^{(k)}\},
\]
defined when the substitution is defined; it need not respect classes. In
particular a parameter of $\theta$ that is a row variable accepts only a record
type.

\section{Type Schemes}\label{sec:static-core-schemes}

A type scheme $\sigma = \forall\tyrowvar^{(k)}.\tau$ \emph{generalises} a type
$\tau'$, written $\sigma \succ \tau'$, if $\tau' =
\tau\{\tau^{(k)}/\tyrowvar^{(k)}\}$ for some $\tau^{(k)}$ for which the
substitution is defined and respects classes; $\tau'$ is then an
\emph{instance} of $\sigma$. If $\sigma' = \forall\tyrowvar'^{(l)}.\tau'$, then
$\sigma$ generalises $\sigma'$, written $\sigma \succ \sigma'$, if $\sigma \succ
\tau'$ and no $\tyrowvar'_j$ occurs free in $\sigma$. By the property of
Section~\ref{sec:static-core-types}, if $\sigma \succ \sigma'$ and $\sigma' \succ
\tau''$ then $\sigma \succ \tau''$.

Two type schemes are equal if they can be obtained from each other by renaming
and reordering bound variables, and by deleting bound variables that do not
occur in the body, other than the variables of class $\knum$ that closure adds
to the scheme of a binding of a group (Section~\ref{sec:static-core-closure}),
which are kept. % A93 In renaming, a type variable is replaced by a type variable
of the same class and a row variable by a row variable. We consider a type
$\tau$ to be a type scheme, identifying it with $\forall().\tau$. For a finite
set $U$ of type and row variables, $\forall U.\tau$ means
$\forall\tyrowvar^{(k)}.\tau$ where $\tyrowvar^{(k)}$ is any enumeration of $U$.

With each special constant $\scon$ we associate a type scheme $\typeof(\scon)$:
it is $\forall\alpha.\alpha$, where $\clsof(\alpha) = \knum$, if $\scon$ is an
integer constant, and the type $\tFloat$, $\tChar$ or $\tString$ if $\scon$ is
a floating point, character or string constant. Thus the instances of
$\typeof(\scon)$ for an integer constant are the types in $\Adm{\knum}$.

\section{Scope of Explicit Type Variables}\label{sec:static-core-explicit}

The type variables written in a type annotation are \emph{explicit}: within the
annotated definition each stands for one type or row variable that the
definition may not instantiate. Explicit variables in scope are recorded in $U
\of C$.

For a set $U$ of type and row variables, $\names U$ is the set of type
variables $\tyvar$ of the syntax such that $\tyvar = \tyrowvar$ for some
$\tyrowvar \in U$. The set $U$ is \emph{unambiguous} if no $\tyvar$ names two
members of $U$. We consider only contexts $C$ in which $U \of C$ is
unambiguous; the rules preserve this.

Let $\ty$ be the type expression of an annotation elaborated in context $C$. A
set $U'$ of type and row variables is \emph{new for $\ty$ in $C$} if
\[
\begin{array}{@{}l@{}}
  \names U' = (\tyvars \ty) \setminus \names(U \of C)\ \text{and}\\
  (U \of C) \cup U'\ \text{is unambiguous}.
\end{array}
\]
Thus $U'$ holds, for each type variable of $\ty$ that is not already explicit
in $C$, the type variable or the row variable that it names; which of the two
is fixed by the rules for type expressions, which accept only variables of $U
\of C$. The members of $U'$ are added to $U$ while the annotation and the body of
its definition elaborate, so that an annotation nested within the body that
writes one of these names means the same variable. % A21
They are bound in the type scheme of the annotated variable, and they must not
occur free in the context of the definition: an explicit variable does not
escape its annotation. % A20

\section{Closure}\label{sec:static-core-closure}

Let $A$ be a semantic object and $\sigma = \forall\tyrowvar^{(k)}.\tau$ a type
scheme. Then $\Clos_A \sigma$, the \emph{closure} of $\sigma$ with respect to
$A$, is the type scheme $\forall\tyrowvar^{(k)}\tyrowvar'^{(l)}.\tau$, where
$\tyrowvar'^{(l)}$ enumerates $\tyvars\sigma \setminus \tyvars A$. In
particular, for a type $\tau$, $\Clos_A \tau = \forall\tyrowvar^{(l)}.\tau$
where $\tyrowvar^{(l)}$ enumerates $\tyvars\tau \setminus \tyvars A$. For a
value environment, $\Clos_A \VE$ is the value environment obtained from $\VE$
by replacing each pair $(\sigma, \is)$ in its range by $(\Clos_A \sigma, \is)$.

Commonly $A$ is a context $C$. Since $U \of C \subseteq \tyvars C$, closure
never generalises an explicit variable in scope. Closure applies to every group
of value bindings, pattern bindings included, whatever the form of their
expressions.

A group of value bindings is closed as one. When
rules~\ruleref{r:s-dec-group} and~\ruleref{r:sm-topdec-value} take
$\Clos_C \VE$ of the value environment $\VE$ of a group, the scheme $\sigma$ of
each binding is closed also over the type variables of class $\knum$ in
$\tyvars\VE \setminus (\tyvars\sigma \cup \tyvars C)$: those of the other
bindings of the group that its own type does not mention. Since such a variable
does not occur in the binding's type, it changes no instance of the scheme. It
matters only to number specialisation
(Section~\ref{sec:dynamic-core-numbers}): an occurrence of the binding gives it
a type that nothing fixes, which is then $\tFloat$. % A93

\section{Inference Rules}\label{sec:static-core-rules}

Each rule of the semantics allows inferences among sentences of the form
\[
  \jdg{A}{\mathit{phrase}}{A'}
\]
where $A$ is usually a context, $\mathit{phrase}$ is a phrase of the Core, and
$A'$ is a semantic object, usually a type, a row, a value environment, or a pair
of a value environment and a type. It is pronounced ``$\mathit{phrase}$
elaborates to $A'$ in (context) $A$''. Some rules have extra hypotheses not of
this form; they are called \emph{side conditions}.

In the presentation of the rules, phrases within single angle brackets
$\opt{\ }$ are called \emph{first options}, and those within double angle
brackets $\optt{\ }$ are called \emph{second options}. To reduce the number of
rules, we have adopted the following convention:
\begin{quote}
In each instance of a rule, the first options must be either all present or all
absent; similarly the second options must be either all present or all absent.
\end{quote}

No Core phrase generates type names, so $T \of C$ is the same in every sentence
of a derivation for a Core phrase; it is used by the static semantics of
Modules (Chapter~\ref{ch:static-modules}). A phrase elaborates to $A'$ when some
derivation of the sentence exists; where the premises of a rule leave a type in
its conclusion undetermined, the derivation may choose any type that satisfies
them. The relation $\vdash$ is overloaded for all syntactic classes.

\rulegroup{Atomic Expressions}{\jdg{C}{\atexp}{\tau}}

\defnrule[r:s-atexp-scon]{\typeof(\scon) \succ \tau}{\jdg{C}{\scon}{\tau}}

\defnrule[r:s-atexp-longvar]{C(\longvar) = (\sigma, \is) \premsep \sigma \succ \tau}{\jdg{C}{\longvar}{\tau}}

\defnrule[r:s-atexp-longcon]{C(\longcon) = (\sigma, \is) \premsep \sigma \succ \tau}{\jdg{C}{\longcon}{\tau}}

\defnrule[r:s-atexp-op]{\OTstat(\opv) = (\sigma, \is) \premsep \sigma \succ \tau}{\jdg{C}{\kw{(}\opv\kw{)}}{\tau}}

\defnrule[r:s-atexp-unit]{}{\jdg{C}{\kw{()}}{\tunit}}

\defnrule[r:s-atexp-tuple]{\jdg{C}{\expm_i}{\tau_i} \quad (1 \le i \le k)}{\jdg{C}{\kw{(}\,\expm_1\kw{,}\ \cdots\kw{,}\ \expm_k\,\kw{)}}{(\tau_1, \cdots, \tau_k)\tTuple{k}}}

\defnrule[r:s-atexp-nil]{}{\jdg{C}{\kw{[}\ \kw{]}}{\tau\,\tList}}

\defnrule[r:s-atexp-record]{\opt{\jdg{C}{\exprow}{\varrho}}}{\jdg{C}{\kw{\kwlb}\,\opt{\exprow}\,\kw{\kwrb}}{\rec{\emptymap\opt{{}+\varrho}}}}

\defnrule[r:s-atexp-update]{\jdg{C}{\var}{\tau} \premsep \tau \in \RecType \premsep \jdg{C}{\exprow}{\varrho} \premsep \varrho \subseteq (\varrho \of \tau)}{\jdg{C}{\kw{\kwlb}\ \var\ \kw{|}\ \exprow\ \kw{\kwrb}}{\tau}}

\defnrule[r:s-atexp-let]{\jdg{C}{\dec}{\VE} \premsep \jdg{C + \VE}{\expm}{\tau}}{\jdg{C}{\kw{let}\ \dec\ \kw{in}\ \expm}{\tau}}

\defnrule[r:s-atexp-paren]{\jdg{C}{\expm}{\tau}}{\jdg{C}{\kw{(}\,\expm\,\kw{)}}{\tau}}

\begin{rulecomments}
\rcomment{r:s-atexp-scon} For an integer constant, $\tau$ may be any type in
$\Adm{\knum}$: $\tInt$, $\tFloat$ or a type variable of class $\knum$. A
floating point, character or string constant has exactly its type. The type
that counts for evaluation is the one the principal derivation gives
(Section~\ref{sec:static-core-principal}); where it is a variable that nothing
binds, number specialisation takes it to be $\tFloat$
(Section~\ref{sec:dynamic-core-numbers}).

\rcomment{r:s-atexp-longvar} The instantiation of type schemes allows different
occurrences of one $\longvar$ to take different types. The identifier status is
not used in this rule, nor in rules \ruleref{r:s-atexp-longcon} and
\ruleref{r:s-atexp-op}; in particular the constructor of a record type alias,
of status $\isv$, elaborates by rule \ruleref{r:s-atexp-longcon} as any
constructor does.

\rcomment{r:s-atexp-op} Infixed applications of operators are derived forms of
this phrase (Appendix~\ref{ch:appendix-derived-forms}), so this rule gives the
type of every application of an operator. The context $C$ is absent from the
premise, so an operator has the same type scheme in every module, whatever the
module declares or imports: a declaration of \kw{add} in a program does not
change the meaning of \kw{+}.

\rcomment{r:s-atexp-tuple} By the grammar $k$ is 2 or 3.

\rcomment{r:s-atexp-nil} Note that $\tau$ does not occur in any premise; thus
\kw{[ ]} has type $\tau\,\tList$ for every $\tau$.

\rcomment{r:s-atexp-record} A record expression has a closed record type: its
fields are exactly those written.

\rcomment{r:s-atexp-update} The side conditions require the type of $\var$ to be
a record type that has every updated field, each at the type of its new
value. The result has the type of $\var$, so an update neither adds a field
nor changes the type of one: if $\kw{r}$ has type $\rec{\{\kw{x} \mapsto
\tInt\}}$ then \kw{\kwlb\ r | x = "s" \kwrb} does not elaborate. The type of
$\var$ may be open, so a function may update a field of any record that has it.

\rcomment{r:s-atexp-let} The variables bound by $\dec$ are in scope in $\expm$
at the type schemes that rule \ruleref{r:s-dec-group} gives them, so $\expm$ may
use them polymorphically.
\end{rulecomments}

\rulegroup{Expression Rows}{\jdg{C}{\exprow}{\varrho}}

\defnrule[r:s-exprow]{\jdg{C}{\expm}{\tau} \premsep \opt{\jdg{C}{\exprow}{\varrho}}}{\jdg{C}{\lab\ \kw{=}\ \expm\ \opt{\kw{,}\ \exprow}}{\{\lab \mapsto \tau\}\opt{{}+\varrho}}}

\begin{rulecomments}
\rcomment{r:s-exprow} The syntactic restrictions ensure $\lab \notin
\Dom\varrho$.
\end{rulecomments}

\rulegroup{Expressions}{\jdg{C}{\expm}{\tau}}

\defnrule[r:s-exp-atexp]{\jdg{C}{\atexp}{\tau}}{\jdg{C}{\atexp}{\tau}}

\defnrule[r:s-exp-app]{\jdg{C}{\expm}{\tau' \rightarrow \tau} \premsep \jdg{C}{\atexp}{\tau'}}{\jdg{C}{\expm\ \atexp}{\tau}}

\defnrule[r:s-exp-lambda]{\jdg{C}{\atpat}{(\VE, \tau)} \premsep \jdg{C + \VE}{\expm}{\tau'}}{\jdg{C}{\kw{\textbackslash}\ \atpat\ \kw{->}\ \expm}{\tau \rightarrow \tau'}}

\defnrule[r:s-exp-case]{\jdg{C}{\expm}{\tau} \premsep \jdg{C}{\match}{\tau \rightarrow \tau'}}{\jdg{C}{\kw{case}\ \expm\ \kw{of}\ \match}{\tau'}}

\begin{rulecomments}
\rcomment{r:s-exp-atexp} The relational symbol $\vdash$ is overloaded for all
syntactic classes (here atomic expressions and expressions).

\rcomment{r:s-exp-lambda} The variables bound by $\atpat$ receive types, not
type schemes, so a parameter has one type throughout $\expm$. Restriction~2 of
Section~\ref{sec:static-core-further} requires $\atpat$ to be exhaustive.
\end{rulecomments}

\rulegroup{Matches}{\jdg{C}{\match}{\tau}}

\defnrule[r:s-match]{\jdg{C}{\mrule}{\tau} \premsep \opt{\jdg{C}{\match}{\tau}}}{\jdg{C}{\mrule\ \opt{\kw{;}\ \match}}{\tau}}

\rulegroup{Match Rules}{\jdg{C}{\mrule}{\tau}}

\defnrule[r:s-mrule]{\jdg{C}{\pat}{(\VE, \tau)} \premsep \jdg{C + \VE}{\expm}{\tau'}}{\jdg{C}{\pat\ \kw{->}\ \expm}{\tau \rightarrow \tau'}}

\begin{rulecomments}
\rcomment{r:s-mrule} This rule lets new free type variables enter the context:
the types of the variables of $\pat$ are chosen in elaborating $\pat$, and they
may have to agree with the explicit variables of annotations within $\expm$.
\end{rulecomments}

\rulegroup{Declarations}{\jdg{C}{\dec}{\VE}}

\defnrule[r:s-dec-group]{\jdg{C + \VE}{\valbind}{\VE} \premsep \VE' = \Clos_C \VE}{\jdg{C}{\valbind}{\VE'}}

\begin{rulecomments}
\rcomment{r:s-dec-group} A declaration is one group of mutually dependent
bindings; a \kw{let} with several definitions is a derived form that divides
them into groups (Appendix~\ref{ch:appendix-derived-forms}). Modifying $C$ by
$\VE$ on the left puts every member of the group in scope in every member's
expression. Since $\VE$ is the same on both sides, an unannotated member has one
type throughout the group: it is monomorphic within the group, and is
generalised only by the closure, for the body of the \kw{let}. % A24
An annotated member enters $\VE$ with the type scheme of its annotation, so it
may be used polymorphically within the group, in its own body included.

The closure applies to every member whatever its expression, so there is no
value restriction: in \kw{let n = 5 in ( n // 2, n / 2 )} the variable
\kw{n} is used at $\tInt$ and at $\tFloat$. The closure leaves the scheme of an
annotated member unchanged, since its free variables are explicit variables of
$C$.
\end{rulecomments}

\rulegroup{Value Bindings}{\jdg{C}{\valbind}{\VE}}

\defnrule[r:s-valbind-var]{\opt{U' \text{ is new for } \ty \text{ in } C \premsep \jdg{C + U'}{\ty}{\tau} \premsep U' \cap \tyvars C = \emptyset} \\ \jdg{C\opt{{}+U'}}{\expm}{\tau} \premsep \optt{\jdg{C}{\valbind}{\VE}}}{\jdg{C}{\var\ \opt{\kw{:}\ \ty}\ \kw{=}\ \expm\ \optt{\kw{and}\ \valbind}}{\{\var \mapsto (\opt{\forall U'.}\tau, \isv)\}\ \optt{{}+\VE}}}

\defnrule[r:s-valbind-pat]{\jdg{C}{\atpat}{(\VE, \tau)} \premsep \jdg{C}{\expm}{\tau}}{\jdg{C}{\atpat\ \kw{=}\ \expm}{\VE}}

\begin{rulecomments}
\rcomment{r:s-valbind-var} When the second option is present,
$\var \notin \Dom\VE$ by the syntactic restrictions. With an annotation, $U'$ holds the variables of $\ty$
not already explicit in $C$ (Section~\ref{sec:static-core-explicit}). They are
in $U$ while $\ty$ and $\expm$ elaborate, so an annotation nested in $\expm$
that writes one of their names means the same variable. % A21
Since $\tau$ is fixed by $\ty$, $\expm$ must have exactly that type, with the
variables of $U'$ uninstantiated: an annotation may be less general than the
type $\expm$ could have, but never more general. Thus \kw{f : a -> a} with
\kw{f x = x + 1} does not elaborate, nor does \kw{g : number -> number} with
\kw{g x = x // 2}.

The condition $U' \cap \tyvars C = \emptyset$ prevents an explicit variable
from escaping its annotation. % A20
In \kw{f x = let g : a -> a ; g y = x in g x}, the body of \kw{g} has type
\kw{a -> a} only if \kw{x} has the explicit type \kw{a}, which is then free in
the context of \kw{g}; so the expression does not elaborate.

\rcomment{r:s-valbind-pat} A pattern binding is a group by itself. The variables
of $\atpat$ receive types here and are generalised by rule
\ruleref{r:s-dec-group}. That rule also puts them in scope in $\expm$, but the
syntactic restrictions forbid $\expm$ to use them, since no pattern binding
lies on a cycle of definitions. % A26
A binding $\var\ \kw{=}\ \expm$ is read by rule \ruleref{r:s-valbind-var}.
\end{rulecomments}

\rulegroup{Type Expressions}{\jdg{C}{\ty}{\tau}}

\defnrule[r:s-ty-tyvar]{\tyvar = \alpha \premsep \alpha \in U \of C}{\jdg{C}{\tyvar}{\alpha}}


\defnrule[r:s-ty-rowvar]{\tyvar = \epsilon \premsep \epsilon \in U \of C}{\jdg{C}{\tyvar}{\orec{\emptymap}{\epsilon}}}


\defnrule[r:s-ty-record]{\opt{\jdg{C}{\tyrow}{\varrho}}}{\jdg{C}{\kw{\kwlb}\,\opt{\tyrow}\,\kw{\kwrb}}{\rec{\emptymap\opt{{}+\varrho}}}}


\defnrule[r:s-ty-openrecord]{\tyvar = \epsilon \premsep \epsilon \in U \of C \premsep \jdg{C}{\tyrow}{\varrho}}{\jdg{C}{\kw{\kwlb}\ \tyvar\ \kw{|}\ \tyrow\ \kw{\kwrb}}{\orec{\varrho}{\epsilon}}}


\defnrule[r:s-ty-unit]{}{\jdg{C}{\kw{()}}{\tunit}}


\defnrule[r:s-ty-tuple]{\jdg{C}{\ty_i}{\tau_i} \quad (1 \le i \le k)}{\jdg{C}{\kw{(}\,\ty_1\kw{,}\ \cdots\kw{,}\ \ty_k\,\kw{)}}{(\tau_1, \cdots, \tau_k)\tTuple{k}}}


\defnrule[r:s-ty-con]{C(\longtycon) = (\theta, \VE) \premsep \jdg{C}{\ty_i}{\tau_i} \quad (1 \le i \le k)}{\jdg{C}{\longtycon\ \ty_1 \cdots \ty_k}{\theta(\tau_1, \cdots, \tau_k)}}


\defnrule[r:s-ty-fun]{\jdg{C}{\ty}{\tau} \premsep \jdg{C}{\ty'}{\tau'}}{\jdg{C}{\ty\ \kw{->}\ \ty'}{\tau \rightarrow \tau'}}


\defnrule[r:s-ty-paren]{\jdg{C}{\ty}{\tau}}{\jdg{C}{\kw{(}\,\ty\,\kw{)}}{\tau}}

\begin{rulecomments}
\rcomment{r:s-ty-tyvar} Rules \ruleref{r:s-ty-tyvar} and \ruleref{r:s-ty-rowvar}
both cover the phrase $\tyvar$. Since $U \of C$ is unambiguous, at most one of
them applies, according to whether the variable of $U \of C$ named by $\tyvar$
is a type variable or a row variable. A type variable that is not explicit in
$C$ does not elaborate: every type variable of an annotation is made explicit by
rule \ruleref{r:s-valbind-var}, and Chapter~\ref{ch:static-modules} makes the
parameters of type declarations explicit in the same way.

\rcomment{r:s-ty-rowvar} A row variable written alone denotes the open record
type with no fields of its own: the record of the fields that $\epsilon$ stands
for.

\rcomment{r:s-ty-openrecord} Only a type variable of class $\kany$ names a row
variable, so \kw{\kwlb\ number | x : Int \kwrb} does not elaborate. % A36

\rcomment{r:s-ty-tuple} By the grammar $k$ is 2 or 3.

\rcomment{r:s-ty-con} For $\theta(\tau_1, \cdots, \tau_k)$ to be defined,
$\theta$ must have arity $k$, so a type constructor is never partially applied.
Application of a type function performs the substitution, so a type alias
denotes its expansion, and an alias with a row parameter applies only to a
record type.
\end{rulecomments}

\rulegroup{Type-expression Rows}{\jdg{C}{\tyrow}{\varrho}}

\defnrule[r:s-tyrow]{\jdg{C}{\ty}{\tau} \premsep \opt{\jdg{C}{\tyrow}{\varrho}}}{\jdg{C}{\lab\ \kw{:}\ \ty\ \opt{\kw{,}\ \tyrow}}{\{\lab \mapsto \tau\}\opt{{}+\varrho}}}

\begin{rulecomments}
\rcomment{r:s-tyrow} The syntactic restrictions ensure $\lab \notin
\Dom\varrho$.
\end{rulecomments}

\rulegroup{Atomic Patterns}{\jdg{C}{\atpat}{(\VE, \tau)}}

\defnrule[r:s-atpat-wild]{}{\jdg{C}{\kw{\_}}{(\emptymap, \tau)}}


\defnrule[r:s-atpat-scon]{\typeof(\scon) \succ \tau \premsep \tau \in \{\tInt, \tChar, \tString\}}{\jdg{C}{\scon}{(\emptymap, \tau)}}


\defnrule[r:s-atpat-var]{}{\jdg{C}{\var}{(\{\var \mapsto (\tau, \isv)\}, \tau)}}

\defnrule[r:s-atpat-con]{C(\longcon) = (\sigma, \isc) \premsep \sigma \succ \tau^{(k)}t}{\jdg{C}{\longcon}{(\emptymap, \tau^{(k)}t)}}

\defnrule[r:s-atpat-unit]{}{\jdg{C}{\kw{()}}{(\emptymap, \tunit)}}

\defnrule[r:s-atpat-tuple]{\jdg{C}{\pat_i}{(\VE_i, \tau_i)} \quad (1 \le i \le k)}{\jdg{C}{\kw{(}\,\pat_1\kw{,}\ \cdots\kw{,}\ \pat_k\,\kw{)}}{(\VE_1 + \cdots + \VE_k,\ (\tau_1, \cdots, \tau_k)\tTuple{k})}}

\defnrule[r:s-atpat-nil]{}{\jdg{C}{\kw{[}\ \kw{]}}{(\emptymap, \tau\,\tList)}}

\defnrule[r:s-atpat-record]{\opt{\jdg{C}{\patrow}{(\VE, \varrho)}} \premsep \tau \in \RecType \premsep \emptymap\opt{{}+\varrho} \subseteq (\varrho \of \tau)}{\jdg{C}{\kw{\kwlb}\,\opt{\patrow}\,\kw{\kwrb}}{(\emptymap\opt{{}+\VE},\ \tau)}}

\defnrule[r:s-atpat-paren]{\jdg{C}{\pat}{(\VE, \tau)}}{\jdg{C}{\kw{(}\,\pat\,\kw{)}}{(\VE, \tau)}}

\begin{rulecomments}
\rcomment{r:s-atpat-scon} The side conditions give an integer constant the type
$\tInt$ in a pattern, never $\tFloat$ or a type variable of class $\knum$; a
character or string constant has its own type. Thus \kw{case 1.5 of 1 -> 0 ; \_
-> 1} does not elaborate. Floating point constants are not patterns, by the
grammar.

\rcomment{r:s-atpat-var} The variable receives a type, not a type scheme.
Constructors and variables are spelt differently, so no status test is needed
here.

\rcomment{r:s-atpat-con} Only status $\isc$ is accepted, so the constructor of a
record type alias, of status $\isv$, cannot occur in a pattern. The instance
must be a constructed type, so this rule accepts only a constructor without
arguments; a constructor with arguments is covered by rule
\ruleref{r:s-pat-con}.

\rcomment{r:s-atpat-tuple} By the grammar $k$ is 2 or 3, and by the syntactic
restrictions the domains of the $\VE_i$ are disjoint.

\rcomment{r:s-atpat-nil} As in rule \ruleref{r:s-atexp-nil}, $\tau$ is
arbitrary.

\rcomment{r:s-atpat-record} The record type $\tau$ may have fields besides those
the pattern names, and it may be open: a record pattern matches every record
with the fields it names.
\end{rulecomments}

\rulegroup{Pattern Rows}{\jdg{C}{\patrow}{(\VE, \varrho)}}

\defnrule[r:s-patrow]{\var = \lab \premsep \opt{\jdg{C}{\patrow}{(\VE, \varrho)}}}{\jdg{C}{\lab\ \opt{\kw{,}\ \patrow}}{(\{\var \mapsto (\tau, \isv)\}\opt{{}+\VE},\ \{\lab \mapsto \tau\}\opt{{}+\varrho})}}

\begin{rulecomments}
\rcomment{r:s-patrow} A field of a record pattern binds the variable spelt as its
label. The syntactic restrictions ensure $\lab \notin \Dom\varrho$, and hence
$\var \notin \Dom\VE$.
\end{rulecomments}

\rulegroup{Patterns}{\jdg{C}{\pat}{(\VE, \tau)}}

\defnrule[r:s-pat-atpat]{\jdg{C}{\atpat}{(\VE, \tau)}}{\jdg{C}{\atpat}{(\VE, \tau)}}

\defnrule[r:s-pat-con]{C(\longcon) = (\sigma, \isc) \premsep \sigma \succ \tau_1 \rightarrow \cdots \rightarrow \tau_n \rightarrow \tau'^{(k)}t \\ \jdg{C}{\atpat_i}{(\VE_i, \tau_i)} \quad (1 \le i \le n)}{\jdg{C}{\longcon\ \atpat_1 \cdots \atpat_n}{(\VE_1 + \cdots + \VE_n,\ \tau'^{(k)}t)}}

\defnrule[r:s-pat-cons]{\jdg{C}{\pat_1}{(\VE_1, \tau)} \premsep \jdg{C}{\pat_2}{(\VE_2, \tau\,\tList)}}{\jdg{C}{\pat_1\ \kw{::}\ \pat_2}{(\VE_1 + \VE_2,\ \tau\,\tList)}}

\defnrule[r:s-pat-as]{\jdg{C}{\pat}{(\VE, \tau)}}{\jdg{C}{\pat\ \kw{as}\ \var}{(\VE + \{\var \mapsto (\tau, \isv)\},\ \tau)}}

\begin{rulecomments}
\rcomment{r:s-pat-con} The instance must be a constructed type after exactly $n$
arguments, so a constructor in a pattern is applied to as many patterns as it
takes. Thus a constructed pattern with too few or too many arguments does not
elaborate. The syntactic restrictions ensure that the domains of the $\VE_i$
are disjoint. % audit Y7

\rcomment{r:s-pat-cons} The type name is $\tList$ of the initial basis, whatever
the program binds to \kw{List}.

\rcomment{r:s-pat-as} The syntactic restrictions ensure $\var \notin \Dom\VE$.
\end{rulecomments}

\section{Principal Derivations}\label{sec:static-core-principal}

A phrase may elaborate in more than one way. An integer constant may take any
type in $\Adm{\knum}$ (rule~\ruleref{r:s-atexp-scon}), and an occurrence of a
variable any instance of its type scheme (rule~\ruleref{r:s-atexp-longvar}),
so two derivations for one phrase may give its integer constants and its
occurrences different types. Whether a phrase elaborates does not depend on the
choice, but evaluation does: \kw{2 \textasciicircum{} -1} is $0$ at $\tInt$
and $0.5$ at $\tFloat$. The dynamic semantics is therefore given for one
derivation, the \emph{principal} one, which gives no type that the phrase does
not force; the further restrictions (Section~\ref{sec:static-core-further}) are
judged on it too, although for them the choice makes no difference. % A95, A96

In a derivation for a phrase, each node is a sentence about an occurrence of a
phrase within it; the nodes of two derivations for the same phrase correspond
when they are about the same occurrence. Let $D$ and $D'$ be derivations for
the same phrase in the same context $C_0$, the type and row variables of $D$
that are not free in $C_0$ having been renamed apart from those of $D'$. Then
$D'$ is an \emph{instance} of $D$ when there is a substitution
$S = \{\tau^{(k)}/\tyrowvar^{(k)}\}$ (Section~\ref{sec:static-core-types}) that
respects classes, with no $\tyrowvar_i$ free in $C_0$, defined on every type of
$D$, such that at each pair of corresponding nodes, with contexts $C$ in $D$ and
$C'$ in $D'$,
\begin{itemize}
\item $S$ takes each variable in $U \of C$ to a variable, and $U \of C'$ is the
  set of the variables so obtained;
\item each type in the sentence of $D'$ is $S$ applied to the corresponding
  type in the sentence of $D$; and
\item each type scheme $\sigma'$ in the sentence of $D'$ is no more general
  than $S\sigma$, for the corresponding type scheme $\sigma$ of $D$: every type
  $\tau$ with $\sigma' \succ \tau$ has $S\sigma \succ \tau$.
\end{itemize}
Since $S$ is defined on every type of $D$, it never gives a record type a
second field with a label it already has; no Core phrase generates type names,
so $T \of C = T \of C'$. The third clause allows a derivation to give a binding
a less general scheme than another derivation gives it, as when
\kw{y = 2 \textasciicircum{} 64 * 4} is given the type $\tInt$ where another
derivation gives it the scheme $\forall\alpha.\alpha$ with $\alpha$ of class
$\knum$. A derivation $D$ for a phrase in a context is \emph{principal} when
every derivation for the phrase in that context is an instance of $D$. Chapter
\ref{ch:static-modules} extends the definition to modules
(Section~\ref{sec:static-modules-principal}).

\begin{property}
Every Core phrase that elaborates in a context $C_0$ has a principal derivation
in $C_0$, and any two of its principal derivations are instances of each other,
so that they differ only by a renaming of type and row variables not free in
$C_0$.
\end{property}

This is the principal type property of Damas and Milner, extended to classes,
rows and annotations. Unification that respects classes has a most general
solution whenever it has a solution, since two classes have a greatest common
subclass or none: $\knum$ and $\kcomp$ give $\knum$, $\kapp$ and $\kcomp$ give
$\kcapp$, and $\knum$ and $\kapp$ none (Section~\ref{sec:static-core-classes}).
The labels of a row are distinct (Section~\ref{sec:static-core-types}), the
explicit variables of an annotation are fixed throughout its scope
(Section~\ref{sec:static-core-explicit}), and a group of value bindings is
monomorphic within itself (rule~\ruleref{r:s-dec-group}), so the inference
algorithm of Damas and Milner, extended in these ways, finds a principal
derivation. The proof is not given here. % A23, A24
% A95. Proof planned in Lean: plans/lean-full-definition.md, step 11
% (complete inference).

Only the variables of class $\knum$ of a principal derivation affect
evaluation: the value of an integer constant, and the result of a number
operation of the initial basis, depend on whether a type is $\tInt$ or
$\tFloat$, while no value depends on how a variable of another class, or of
none, is instantiated. Number specialisation
(Section~\ref{sec:dynamic-core-numbers}) starts from the principal derivation
and fixes its variables of class $\knum$.

\section{Further Restrictions}\label{sec:static-core-further}

There are a few restrictions on programs which are better expressed apart from
the inference rules. They are judged on the principal derivation
(Section~\ref{sec:static-core-principal}): the type $\tau$ with which the
patterns of a match elaborate, and the type of a pattern, are those that
derivation gives. A phrase that elaborates by the rules but whose principal
derivation breaks one of them is not legal; it is treated as a phrase that does
not elaborate. In fact every derivation for the phrase gives the same answer.
A pattern that the principal derivation elaborates at a type variable, at the
top of a match or within another pattern, is a wildcard, a variable, or one of
these with \kw{as} or parentheses, since every other pattern (a special
constant, a constructor, \kw{()}, a tuple, \kw{[\ ]}, \kw{::} or a record
pattern) elaborates only at a type that is not a variable; and such a pattern
covers every inhabitant. Another derivation may replace the variable by any
type, which has inhabitants (below), and the same patterns cover them. % A96, A97

The first two restrictions speak of the \emph{inhabitants} of a type, which are
phrases, possibly infinite, built from constructors and special constants, and
of the patterns that \emph{cover} them. The \emph{constructors} of a type name $t$ are those
that the declaration of the custom type with type name $t$ binds with status
$\isc$ (Chapter~\ref{ch:static-modules}), or, for a type name of the initial
basis, those that Appendix~\ref{ch:appendix-static-basis} gives it; the
constructors of $\tList$ are \kw{[]} and \kw{::}. The inhabitants of a type
$\tau$ are given by the first of the following clauses that applies:
\begin{itemize}
\item if $\tau$ is $\tInt$, $\tChar$ or $\tString$, each special constant
  $\scon$ such that $\typeof(\scon) \succ \tau$;
\item if $\tau = \tunit$, the phrase \kw{()};
\item if $\tau = (\tau_1, \cdots, \tau_k)\tTuple{k}$, the phrases
  $\kw{(}h_1\kw{,}\allowbreak\ \cdots\kw{,}\allowbreak\ h_k\kw{)}$, where each
  $h_i$ is an inhabitant of $\tau_i$;
\item if $\tau = \tau'^{(k)}t$ and $t$ has constructors, the phrases
  $\con\ h_1 \cdots h_n$ (written $h_1\ \kw{::}\ h_2$ when $\con$ is
  \kw{::}), where $\con$ is a constructor of $t$ whose type scheme $\sigma$
  satisfies $\sigma \succ \tau_1 \rightarrow \cdots \rightarrow \tau_n
  \rightarrow \tau$, and each $h_i$ is an inhabitant of $\tau_i$;
\item otherwise (a function type, a record type, a type variable, or a type
  whose type name has no constructors), the single phrase \kw{\_}.
\end{itemize}
The clauses are read coinductively: an inhabitant may be infinite, provided
each of its parts is given by a clause. So \kw{type V = V V} has the one
inhabitant \kw{V (V (V \ldots))}, and \kw{Never}, whose one constructor
\kw{JustOneMore} takes a \kw{Never}
(Appendix~\ref{ch:appendix-static-basis}), has the one inhabitant
\kw{JustOneMore (JustOneMore \ldots)}. Every type has an inhabitant, and so
every constructor of a type gives inhabitants of it, at every instance.
Inhabitants are not values: no value has type \kw{Never}. An inhabitant
describes what a pattern may be asked to match, which depends only on the
constructors of a type. Thus with \kw{r} of type \kw{Result Never Int}, the
second rule of \kw{case r of Ok n -> n ; Err e -> never e} covers
\kw{Err (JustOneMore \ldots)} and is not redundant, while with \kw{m} of type
\kw{Maybe Never}, \kw{case m of Nothing -> 0} is not exhaustive. Where every
type that a match involves has a finite inhabitant, the infinite inhabitants
change neither restriction, since a pattern examines only a finite part of an
inhabitant. % A97
The wildcard \kw{\_}, a variable and a record pattern cover every inhabitant.
A special constant covers each special constant that denotes the same
character or string or, for integer constants, has the same value as an
$\tInt$ (Chapter~\ref{ch:syntax-core}), however the two are written. % A31
A pattern $\longcon\ \atpat_1 \cdots \atpat_n$ ($n \geq 0$) covers $\con\ h_1
\cdots h_n$ when $\longcon$ is $\con$ or $\modid.\con$ and each $\atpat_i$
covers $h_i$; likewise \kw{[\ ]} covers \kw{[]}, and $\pat_1\ \kw{::}\ \pat_2$
covers $h_1\ \kw{::}\ h_2$ when $\pat_1$ covers $h_1$ and $\pat_2$ covers
$h_2$. The pattern \kw{()} covers \kw{()}, a tuple pattern covers a tuple
component by component, and $\pat\ \kw{as}\ \var$ and $\kw{(}\pat\kw{)}$ cover
what $\pat$ covers. The restrictions are:

\begin{enumerate}
\item\citelabel{res:core-irredundancy} \emph{Irredundancy.} In a match $\pat_1\ \kw{->}\ \expm_1\ \kw{;}\ \cdots\
  \kw{;}\ \pat_n\ \kw{->}\ \expm_n$ whose patterns elaborate with type $\tau$,
  each $\pat_j$ must cover some inhabitant of $\tau$ that no $\pat_i$ with $i <
  j$ covers. % A31

\item\citelabel{res:core-exhaustiveness} \emph{Exhaustiveness.} In a match as above, every inhabitant of $\tau$
  must be covered by some $\pat_i$. Likewise, in a pattern binding
  $\atpat\ \kw{=}\ \expm$ and in a lambda
  $\kw{\textbackslash}\ \atpat\ \kw{->}\ \expm$, the pattern $\atpat$ must
  cover every inhabitant of its type. The
  restrictions are inherited by derived forms: in a function definition
  $\var\ \atpat_1 \cdots \atpat_n\ \kw{=}\ \expm$, each $\atpat_i$ must be
  exhaustive by itself. In deciding exhaustiveness, no set of special
  constants covers every inhabitant of $\tInt$, $\tChar$ or $\tString$, whatever the range of $\tInt$
  (Appendix~\ref{ch:appendix-implementation}).
\end{enumerate}
